\subsection{完备赋范结合代数中的指数}

在本号中，A 表示带有范数 \(x \mapsto \| x \|\) 的含幺结合代数，并满足以下条件：

\[
\begin{array}{c} \| x + y \| \leqslant \sup (\| x \|, \| y \|) \\ \| x y \| \leqslant \| x \|. \| y \| \\ \| 1 \| = 1 \end{array}
\]

对 A 中的 \(x, y\) 成立，并且关于此范数完备。§ 7 第 3 号的第二、三段结果仍然有效。

取 \(\mathbf{I} = \{\mathrm{U}\}\)，并考察 \(\tilde e\) 中级数 \(\tilde{l}\) 和 \(e(\mathrm{U}) = \sum_{n\geqslant 1}\frac{\mathrm{U}^n}{n!}\) 的像 \(l(\mathrm{U}) = \sum_{n\geqslant 1}(-1)^{n - 1}\frac{U^n}{n}\) 和 \(\hat{\mathrm{P}} (\mathrm{A};\mathrm{A})\)。于是：

\begin{enumerate}
\def\labelenumi{(\arabic{enumi})}
\setcounter{enumi}{19}
\tightlist
\item
\end{enumerate}

\[
\left\| \left(\frac {\mathrm{U} ^ {n}}{n !}\right) ^ {\sim} \right\| \leqslant a ^ {- (n - 1) \theta}\tag{20}
\]

\[
\left\| \left(\frac {\mathrm{U} ^ {n}}{n}\right) ^ {\sim} \right\| \leqslant a ^ {- \frac {\log n}{\log p}}\tag{21}
\]

由第 1 号引理 2。于是级数 \(\tilde{e}\)（分别为 \(\tilde{l}\)）的绝对收敛半径为 \(\geqslant a^\theta\)（分别 \(\geqslant 1\)）（《微分流形与解析流形》R，4.1.3）。对 \(R > 0\)，令 \(G_R = \{x \in A \mid \|x\|<R\}\)；记 \(G = G_\theta\)。级数 \(\tilde{e}\)（分别为 \(\tilde{l}\)）定义了解析映射 \(e_{A}\)（分别为 \(l_{A}\)），其定义域分别为 \(G\) 和 \(G_{1}\)，值域为 A。记：

\begin{enumerate}
\def\labelenumi{(\arabic{enumi})}
\setcounter{enumi}{21}
\tightlist
\item
\end{enumerate}

\[
\exp_ {\mathrm{A}} (x) = 1 + e _ {\mathrm{A}} (x) = \sum_ {n \geqslant 0} \frac {x ^ {n}}{n !}\tag{22}
\]

\[
\log_ {\mathrm{A}} (x) = l _ {\mathrm{A}} (x - 1) = \sum_ {n \geqslant 1} (- 1) ^ {n - 1} \frac {(x - 1) ^ {n}}{n} \quad \text { for } x - 1 \in \dot {\mathrm{G}} _ {1}\tag{23}
\]

（若不会引起混淆，省略下标 A。）对 \(x \in G_R\) 和 \(n \geqslant 1\)，

\[
\left\| \frac {x ^ {n}}{n} \right\| \leqslant \left\| \frac {x ^ {n}}{n !} \right\| <   R ^ {n} a ^ {- (n - 1) \theta} = R \left(\frac {R}{a ^ {\theta}}\right) ^ {n - 1}\tag{24}
\]

从而当 \(e_{\mathrm{A}}(\mathrm{G}_{\mathrm{R}}) \subset \mathrm{G}_{\mathrm{R}}, l_{\mathrm{A}}(\mathrm{G}_{\mathrm{R}}) \subset \mathrm{G}_{\mathrm{R}}\) 时，\(R \leqslant a^{\theta}\)。

命题 4。令 R 为满足 \(0 < R \leqslant a^{\theta}\) 的实数。映射 \(\exp_{\mathbf{A}}\) 定义了从 \(\mathbf{G}_{R}\) 到 \(1 + \mathbf{G}_{R}\) 的解析同构，而逆同构是 \(\log_{\mathbf{A}}\) 在 \(1 + \mathbf{G}_{R}\) 上的限制。

\(e(l(\mathbf{X})) = l(e(\mathbf{X})) = \mathbf{X}\)。由 (20)、(21) 以及《微分流形与解析流形》R，4.1.5，得到 \(e_{\mathrm{A}}(l_{\mathrm{A}}(x)) = l_{\mathrm{A}}(e_{\mathrm{A}}(x))\) 对 \(x \in \mathrm{G}_{\mathrm{R}}\) 成立。于是

\[
\begin{array}{l l} \exp_ {\mathrm{A}} (\log_ {\mathrm{A}} x) = x & \text {for} x \in 1 + \mathrm{G} _ {\mathrm{R}} \\ \log_ {\mathrm{A}} (\exp_ {\mathrm{A}} x) = x & \text {for} x \in \mathrm{G} _ {\mathrm{R}} \end{array}
\]

证明完成。

若在 A 上赋予括号 \([x,y]=xy-yx\)，则 A 成为完备赋范李代数，因为 \(\|xy-yx\|\leqslant\sup(\|xy\|,\|yx\|)\leqslant\|x\|. \|y\|\)。第 3 号命题 2 蕴含 \(\tilde{H}\) 的绝对收敛域包含 \(G\times G\)，因此 \(\tilde{H}\) 定义了解析函数 \(h:G\times G\to A\)；于是

\[
h (x, y) = \sum_ {r, s \geqslant 0} \mathrm{H} _ {r, s} (x, y).\tag{25}
\]

命题 5。对于 G 中的 x、y，

\[
\exp_ {\mathrm{A}} x. \exp_ {\mathrm{A}} y = \exp_ {\mathrm{A}} h (x, y).\tag{26}
\]

由 \(e^{\mathrm{U}}e^{\mathrm{V}} = e^{\mathrm{H(U,V)}}\)，从而

\[
m \circ (1 + \tilde {e}, 1 + \tilde {e}) = (1 + \tilde {e}) \circ \tilde {\mathrm{H}}
\]

在 \(\hat{\mathrm{P}}(\mathbf{A} \times \mathbf{A}; \mathbf{A})\) 中（其中 m 表示 A 上的乘法）。于是由命题 2、引理 3 以及《微分流形与解析流形》R，4.1.5，得到命题。


